% Ableitung und Änderungsrate · Mittlere Änderungsrate · Prüfungs-Fokus Abitur GK – bank/ableitung-und-aenderungsrate/e1.jsonl (zusammenbau v0.6, Prüfungs-Fokus)
\einheitenkopf[e1]{Einheit 1 · Mittlere Änderungsrate und Sekante}
\zweigzeile{Abi GK ×5}
\setcounter{aufgabe}{0}

% Anlauf (ohne Prüfkennung)
\begin{aufgabe}{Mittlere Änderungsrate – Anlauf}
\begin{teile}
\teil „Wie stark hat der Wasserstand in den ersten fünf Stunden durchschnittlich je Stunde zugenommen?“ \\ \kreuz{Zeitraum – Differenzenquotient, Sekante} \\ \kreuz{Zeitpunkt – Ableitung, Tangente}
\teil Um 8 Uhr sind 120 Besucher im Zoo, um 11 Uhr 300. Wie groß ist die mittlere Änderungsrate der Besucherzahl in diesem Zeitraum? \leerfeld Besucher je Stunde
\teil Zwei Modelle beschreiben die Abkühlung eines Getränks: $A(t) = 70 \cdot \mathrm{e}^{-0{,}1t} + 20$ und $B(t) = 90 - 3t$ ($t$ in Minuten, Temperatur in °C). Vergleiche die Beträge der mittleren Änderungsraten beider Modelle im Intervall $[0; 10]$.
\end{teile}
\end{aufgabe}

\begin{aufgabe}{Mittlere Änderungsrate – Prüfungsaufgaben}
\begin{teile}
\teil Der Graph zeigt die Anzahl $a$ der bis zum Zeitpunkt $t$ (Stunden nach 12:00 Uhr) abgegebenen Bewertungen. Lies die Werte für 13:00 Uhr und 15:00 Uhr ab und berechne, wie viele Bewertungen in diesem Zeitraum durchschnittlich je Stunde abgegeben wurden \hfill (A25).

\begin{ksys}[xmin=0,xmax=4,ymin=0,ymax=1800,xstep=1,ystep=200,xlabel=t in h,ylabel=a,ablesen]\funktion{100*\x^2}{a}\end{ksys}
\leerfeld je Stunde
\teil Der Hormonspiegel eines Patienten wird durch $h(t) = 6t \cdot \mathrm{e}^{-0{,}05t} + 40$ beschrieben ($t$ in Tagen seit Behandlungsbeginn, $h(t)$ in Prozent des Sollwerts). Berechne die mittlere Änderungsrate des Hormonspiegels in den ersten fünf Tagen \hfill (A19). \leerfeld Prozentpunkte je Tag
\teil Bei einem Laktattest beschreibt $k(x) = \frac{1}{20} \cdot (x^3 - 24x^2 + 200x - 500)$ für $10 \le x \le 16$ die Laktatkonzentration in mmol/l in Abhängigkeit von der Laufgeschwindigkeit $x$ in km/h. Berechne die mittlere Änderungsrate der Laktatkonzentration zwischen 10 und 16 km/h \hfill (A19). \leerfeld mmol/l je km/h
\teil Das Luftvolumen in der Lunge beim Einatmen wird durch $h(t) = 0{,}4t^2 - 1{,}2t + 2$ beschrieben ($t$ in Sekunden, $h(t)$ in Litern; Minimum bei $t = 1{,}5$). Gib die Bedeutung der Terme I: $h(3) - h(1{,}5)$ und II: $\frac{h(3) - h(1{,}5)}{1{,}5}$ im Sachzusammenhang an und berechne beide \hfill (A26).
\teil Ein Höhenprofil wird für $0 \le x \le 6$ durch $f(x) = (x + 2) \cdot \mathrm{e}^{-0{,}4x}$ beschrieben (1 LE = 1 km). Im Intervall $[1; 5]$ soll das Profil näherungsweise durch die Gerade durch die Punkte $(1 | f(1))$ und $(5 | f(5))$ ersetzt werden. Gib eine Gleichung dieser Geraden an \hfill (A18).
\teil Ein Höhenprofil wird durch $f(x) = (x + 2) \cdot \mathrm{e}^{-0{,}4x}$ beschrieben (1 LE = 1 km). Von einem zweiten Profil $h$ ist bekannt: $h(0) = 1{,}8$ und $h(5) = 0{,}6$. Vergleiche die Beträge der mittleren Steigungen von $f$ und $h$ im Intervall $[0; 5]$ \hfill (A18).
\end{teile}
\end{aufgabe}

\begin{aufgabe}{Mittlere Änderungsrate – Prüfungsaufgaben – weiter}
\begin{teile}
\teil Der Schienenverlauf einer Modelleisenbahn besteht aus vier gleichen Bögen von je $3{,}2$ m Länge und zwei Halbkreisen mit dem Radius $0{,}5$ m, die die Bögen verbinden. Ein Zug fährt im Durchschnitt $4$ m je Minute. Wie lange braucht er für eine volle Runde? \hfill (A26) \\ \leerfeld[min]
\teil Der Bestand eines Rohstoffs wird durch $k(t) = 0{,}1t^3 - 1{,}5t^2 + 100$ beschrieben ($t$ in Jahren, $k(t)$ in Mio. Tonnen). Die Gleichung $\frac{k(c) - k(0)}{c} = -5$ hat genau zwei Lösungen. Bestimme die kleinere Lösung und deute Gleichung und Lösung im Sachzusammenhang \hfill (A26).
\teil Die Kosten für die Produktion von $x$ Kubikmetern einer Flüssigkeit werden für $0 \le x \le 8$ durch $K(x) = x^3 - 9x^2 + 30x + 10$ beschrieben ($K(x)$ in 1\,000 Euro). Aussage: „Je größer die Produktionsmenge, desto höher sind die Kosten für einen zusätzlichen Kubikmeter.“ Beurteile die Aussage \hfill (A18).
\teil Die Kosten für die Produktion von $x$ Kubikmetern einer Flüssigkeit werden für $0 \le x \le 8$ durch $K(x) = 0{,}5x^3 - 6x^2 + 30x + 15$ beschrieben ($K(x)$ in 1\,000 Euro). Aussage: „Je größer die Produktionsmenge, desto höher sind die Kosten für einen zusätzlichen Kubikmeter.“ Beurteile die Aussage \hfill (A18).
\teil Die CO$_2$-Konzentration in einem Raum wird durch $g(x) = -500 \cdot \mathrm{e}^{-0{,}4x} + 900$ beschrieben ($x$ in Stunden, $g(x)$ in ppm). Weise nach, dass die mittlere Änderungsrate über jede Stunde $[x; x + 1]$ durch $500 \cdot (1 - \mathrm{e}^{-0{,}4}) \cdot \mathrm{e}^{-0{,}4x}$ gegeben ist, und berechne auf eine Minute genau, ab wann sie erstmals $80$ ppm je Stunde unterschreitet \hfill (A19).
\teil Die CO$_2$-Konzentration in einem Raum wird durch $g(x) = -400 \cdot \mathrm{e}^{-0{,}3x} + 800$ beschrieben ($x$ in Stunden, $g(x)$ in ppm). Weise nach, dass die mittlere Änderungsrate über jede Stunde $[x; x + 1]$ durch $400 \cdot (1 - \mathrm{e}^{-0{,}3}) \cdot \mathrm{e}^{-0{,}3x}$ gegeben ist, und berechne auf eine Minute genau, ab wann sie erstmals $60$ ppm je Stunde unterschreitet \hfill (A19).
\end{teile}
\end{aufgabe}

\medskip
%% Merkkasten Einheit 1: 7 Zeilen, Vorgabe höchstens fünf (3.1) – ungekürzt
\uebersichtskasten{Differenzenquotient: Die mittlere Änderungsrate von f auf dem Intervall [a; b] ist (f(b) − f(a))/(b − a) – die Änderung von f geteilt durch die Länge des Intervalls. Das ist die Steigung der Sekante durch die Punkte (a | f(a)) und (b | f(b)). \\ f(x) = x² auf [1; 3]: (f(3) − f(1))/(3 − 1) = (9 − 1)/2 = 4. \\ Aus Tabelle oder Graph: die beiden Werte ablesen, die Differenz durch die Länge des Zeitraums teilen – Uhrzeiten vorher in Stunden umrechnen. Einheit: Einheit von f je Einheit von x. \\ Bestand 1200 um 8:00 Uhr und 1500 um 9:30 Uhr: (1500 − 1200)/1,5 = 200 je Stunde. \\ Sekantengleichung: Steigung m als Differenzenquotient, dann y = m · (x − a) + f(a). \\ f(x) = x², a = 1, b = 3: m = 4, y = 4 · (x − 1) + 1 = 4x − 3. \\ Differenz gegen Differenzenquotient: f(b) − f(a) ist die absolute Änderung (Einheit von f), (f(b) − f(a))/(b − a) die mittlere Änderung je Einheit von x. Eine Durchschnittsgeschwindigkeit ist Weg je Zeit – die Zeit ist Weg geteilt durch Geschwindigkeit.}
